\section{Growth of cloud computing}
\label{sec:Growth}
The paradigm of cloud computing presents the fact that it is fulfilling a long-held dream in the information technology industry: computing as a utility \cite{Armbrust:EECS-2009-28}. This transformation has reshaped the way IT hardware is designed, purchased, and utilized, reconstructing the industry from a large up-front capital to a pay-as-you-go economic model. Historically, the growth of cloud computing was observed by the development of large-scale commodity-computer data centers, allowing providers to achieve significant economies of scale in electricity, hardware, and operations \cite{Armbrust:EECS-2009-28}.
 
In addition, \cite{Armbrust:EECS-2009-28} shows three new aspects that mark the initial emergence of cloud adoption. Those are the illusion of infinite computing resources available on demand, the elimination of up-front commitments by users, and the ability to pay for resources on a short-term, as-needed basis. Thus, these aspects give cloud computing elasticity, which allowed organizations to attach resources to workloads much more closely, transferring the risks of overprovisioning and underprovisioning from the service operator to the cloud vendor \cite{Armbrust:EECS-2009-28}. A notable early example of this growth was the Animoto service, which scaled from 50 to 3,500 servers in three days after it was deployed on Facebook.

Furthermore, the growth of cloud computing is characterized by a decentralized infrastructure through emerging architectures such as serverless computing and edge computing \cite{varghese2022cloud}. Serverless architectures (Function-as-a-Service) allow developers to concentrate solely on application logic while the platform manages all underlying scaling, fault tolerance, and backend management. Simultaneously, the increase in Internet of Things (IoT) devices has driven the adoption of edge and fog computing, bringing processing power closer to data sources to meet ultra-low-latency requirements and reduce network bottlenecks \cite{varghese2022cloud}. This continuous expansion into distributed, heterogeneous, and ad hoc environments highlights the increasing complexity of the modern cloud ecosystem.

As technology developed, the landscape evolved from single-provider environments to more complex infrastructures \cite{varghese2022cloud}. Today, hybrid cloud and multi-cloud models have become the dominant standard. \cite{varghese2022cloud}, \cite{flexera2026cloud} indicate that approximately 73\% of organizations now employ a hybrid cloud approach, combining public and private environments to balance cost, performance, and governance. Large enterprises, in particular, have observed a significant increase in cloud spending, with more than 76\% of large organizations spending more than \$5 million monthly on public cloud services by 2026. However, as these hybrid cloud architectures grow in sophistication and monthly spend, they require centralized oversight through a Cloud Center of Excellence (CCOE) or FinOps teams to manage the growing complexity and associated security risks \cite{flexera2026cloud}. Consequently, operational and security processes are required for hybrid cloud infrastructure management. 
